\documentclass[11pt]{article}
\usepackage{amsmath, amssymb, amsthm}
\usepackage{geometry}
\usepackage{hyperref}
\usepackage{booktabs}
\usepackage{enumitem}
\usepackage[dvipsnames]{xcolor}
\usepackage{tcolorbox}
\usepackage{listings}

\geometry{margin=1.05in}
\hypersetup{colorlinks=true, linkcolor=blue!60!black, urlcolor=blue!60!black}

\newcommand{\eps}{\varepsilon}
\newcommand{\code}[1]{\texttt{#1}}
\newcommand{\R}{\mathbb{R}}
\newcommand{\dv}{\nabla\cdot}
\newcommand{\Nc}{\mathcal{N}}
\setlist[itemize]{leftmargin=1.4em, itemsep=2pt}
\setlist[enumerate]{leftmargin=1.8em, itemsep=3pt}

\lstset{basicstyle=\ttfamily\small, columns=fullflexible, keepspaces=true,
        frame=single, framerule=0.3pt, rulecolor=\color{gray!60}, xleftmargin=4pt,
        breaklines=true, language=Python, showstringspaces=false}

\newcounter{ex}
\newenvironment{ex}[1]{\refstepcounter{ex}\par\medskip\noindent\textbf{Exercise \theex}\ \ \textit{#1}\par\nopagebreak\smallskip}{\par\medskip}
\newtcolorbox{codebox}[1]{colback=green!4!white, colframe=green!40!black, title=Code check: \code{#1}, fonttitle=\bfseries\small, boxrule=0.5pt, arc=1pt, left=4pt, right=4pt, top=3pt, bottom=3pt}
\newtcolorbox{goal}{colback=orange!6!white, colframe=orange!70!black, boxrule=0.5pt, arc=1pt, left=4pt, right=4pt, top=3pt, bottom=3pt}
\newtcolorbox{physbox}[1]{colback=blue!4!white, colframe=blue!50!black, title=#1, fonttitle=\bfseries\small, boxrule=0.5pt, arc=1pt, left=4pt, right=4pt, top=3pt, bottom=3pt}

\title{Worksheet: the unbalanced solver\\ and the glymphatic DCE-MRI problem}
\author{Companion to \code{python/optimal\_transport/*\_unbalanced.py},\\
\code{docs/unbalanced\_1d.md}, and Chen, Benveniste, Tannenbaum, Sci.\ Rep.\ 14 (2024)}
\date{24 October 2026}

\begin{document}
\maketitle

\begin{goal}
\textbf{Goal for today.} Be able to (i) write the unbalanced problem and say what $\delta$
does, (ii) derive both LADMM subproblems, the prox and the projection, and see that each is
a one-line change to the balanced code, (iii) say what every symbol means in a rat-brain
tracer scan and translate urOMT's parameters into ours, (iv) build $\rho_0,\rho_1$ from the
public data, and (v) know the three edits that port the 1D solver to 3D. Parts A--C are
pen, paper, and the 1D numpy code (runs on the Mac). Parts D--F are the physics, the data
recipe, and the GPU plan. Answers in Appendix~\ref{sec:answers}.
\end{goal}

\section*{Conventions}
Balanced solver: $\min\iint|m|^2/2\rho$ s.t.\ $\partial_t\rho+\dv m=\eps\Delta\rho$,
$\rho(0)=\rho_0$, $\rho(1)=\rho_1$, $m\cdot n=0$ on the walls of $[0,1]^d$, $t\in[0,1]$.
LADMM splitting (\code{pipeline.py}): $x=(\rho,m)$ on the staggered grid, $y=(\rho,m)$ on
cell centres, constraint $Ax-y=0$ with $A$ the staggered-to-centre interpolation;
$f_1(x)=\iota_{\rm FP}(x)$ handled by a projection, $f_2(y)=$ KE handled by a pointwise prox.
Chen et al.\ write $\rho|v|^2$ with no $\tfrac12$ and $\alpha\chi\rho r^2$ for the source;
we keep our $\tfrac12$ throughout and translate in Part~D.

\section*{Part A. The unbalanced formulation}

\begin{ex}{Write the problem.}
Add a source $\zeta(t,x)$ to the continuity equation and charge for it with the
Fisher--Rao cost. Write out
\[
  \min_{\rho,m,\zeta}\ \iint\frac{|m|^2}{2\rho}+\frac{\delta^2}{2}\frac{\zeta^2}{\rho}\,dx\,dt
  \quad\text{s.t.}\quad \partial_t\rho+\dv m-\eps\Delta\rho=\zeta,\ \ \rho(0)=\rho_0,\ \rho(1)=\rho_1,
\]
with $\int\rho_0\neq\int\rho_1$ allowed. Introduce the growth rate $g=\zeta/\rho$ and the
velocity $v=m/\rho$ and rewrite the integrand as $\tfrac12\rho\,(|v|^2+\delta^2g^2)$. What
are the units of $\delta$? Why must $\delta$ have those units for the two terms to be
comparable?
\end{ex}

\begin{ex}{Pure reaction has a closed form.}
Suppose $\rho_1=c\,\rho_0$ for a constant $c>0$ and $\eps=0$. Guess that the optimum has
$m\equiv0$ and find the path. With $m=0$ the constraint is $\partial_t\rho=\zeta$, so the cost
is $\tfrac{\delta^2}{2}\iint(\partial_t\rho)^2/\rho$. Substitute $\rho=\sigma^2$ and minimize
$2\delta^2\iint(\partial_t\sigma)^2$ pointwise in $x$: show $\sigma$ is linear in $t$, so
\[
  \rho_t=\big((1-t)\sqrt{\rho_0}+t\sqrt{\rho_1}\big)^2,\qquad
  \text{cost}=2\delta^2\!\int\!\big(\sqrt{\rho_1}-\sqrt{\rho_0}\big)^2dx=2\delta^2(\sqrt c-1)^2\!\int\!\rho_0 .
\]
This is the Hellinger geodesic. It is test T1 of \code{unbalanced\_1d.md}: for uniform
$\to2\times$uniform on $[0,1]$ the cost is $2\delta^2(\sqrt2-1)^2$, which with $\delta=0.3$ is
$0.03088$; the solver hit it to all printed digits in 53 iterations. Why does adding $\eps>0$
not change this answer when $\rho_0$ is uniform?
\end{ex}

\begin{ex}{What $\delta$ decides: move or teleport.}
Two unit-mass lumps a distance $d$ apart, $\eps=0$. Compare the cost of (a) transporting the
mass, $\tfrac12d^2$, with (b) destroying it in place and recreating it at the target, which
by Exercise~2 costs $2\delta^2(0-1)^2+2\delta^2(1-0)^2=4\delta^2$. Show the crossover is at
$d\propto\delta$ and that the true WFR geodesic (which shrinks the mass while moving it) is
cheaper than both, so the exact crossover is a pure number times $\delta$; the hard limit
beyond which no transporting path competes is $\pi\delta$. The 1D tests anchor this: T3 with
gap $0.19>\pi\delta=0.157$ teleports (kinetic energy $5\cdot10^{-6}$), T3$'$ with
$\pi\delta=0.31>0.19$ transports (kinetic energy $0.0126$, mass dips to $0.56$ mid-way).
Explain the mass dip in T3$'$ in words.
\end{ex}

\begin{ex}{Optimality conditions.}
Introduce a multiplier $\phi(t,x)$ for the constraint and derive, pointwise where $\rho>0$,
\[
  v=\nabla\phi,\qquad \delta^2g=\phi\quad\Longrightarrow\quad m=\rho\nabla\phi,\ \ \zeta=\rho\phi/\delta^2 .
\]
Conclude the consistency relation $v=\nabla(\delta^2\zeta/\rho)$ between the velocity and the
growth rate, which is test T5 (median relative error $8\cdot10^{-5}$). What happens to this
relation as $\delta\to\infty$, and why does that recover the balanced problem (test T4:
difference to the balanced solver $\propto1/\delta^2$)?
\end{ex}

\section*{Part B. The two subproblems}

The code's trick is $s:=\delta\zeta$. Then the reaction cost is $s^2/2\rho$ and the
constraint is $\partial_t\rho+\dv m-\eps\Delta\rho-s/\delta=0$.

\begin{ex}{The prox is the old prox.}
The $y$-subproblem is, pointwise at each cell,
\[
  \min_{\rho,m,s}\ \frac{|m|^2+s^2}{2\rho}+\frac1{2\sigma}\big(|\rho-\bar\rho|^2+|m-\bar m|^2+|s-\bar s|^2\big),
  \qquad \sigma=1/\gamma .
\]
(a) Differentiate in $m$ and in $s$ and show both shrink by the same factor,
$m=\tfrac{\rho}{\rho+\sigma}\bar m$, $s=\tfrac{\rho}{\rho+\sigma}\bar s$.
(b) Substitute back, differentiate in $\rho$, and show $\rho$ solves the cubic
\[
  \rho^3+(2\sigma-\bar\rho)\rho^2+(\sigma^2-2\sigma\bar\rho)\rho-\sigma\big(\sigma\bar\rho+\tfrac12q^2\big)=0,
  \qquad q^2=|\bar m|^2+\bar s^2 .
\]
(c) Compare with the balanced cubic in \code{prox.py}: the only difference is $q^2$. Why is
the function $(\rho,m,s)\mapsto(|m|^2+s^2)/2\rho$ convex, and what is it the perspective of?
(d) Which cost would \emph{not} reduce to this cubic? Give one example (an entropy-type or
$L^1$ growth penalty) and say what the pointwise solve would become.
\end{ex}

\begin{codebox}{prox\_unbalanced.py}
\code{q2 = mx\_c**2 + s\_c**2}; coefficients \code{b = 2*sigma - rho\_c},
\code{c = sigma**2 - 2*sigma*rho\_c}, \code{d = -sigma*(sigma*rho\_c + 0.5*q2)};
\code{fac = rho\_new/(rho\_new + sigma)}; \code{s\_new = fac*s\_c}. Match (a), (b). The
balanced \code{solve\_cubic} is reused unchanged. In 3D, \code{q2} gains \code{my**2 + mz**2}.
\end{codebox}

\begin{ex}{The projection is the old projection plus a constant.}
The $x$-subproblem is the Euclidean projection of $(\bar\rho,\bar m,\bar s)$ onto the affine
set $\{M x=0\}$, $M=[\,A_{\rm FP},\ -I/\delta\,]$, where $A_{\rm FP}(\rho,m)$ is the discrete
$\partial_t\rho+\dv m-\eps\Delta\rho$ (with $\rho_0,\rho_1$ baked into the time
differences) and the $-I/\delta$ block acts on $s$.
(a) Write the projection as $x=\bar x-M^{\!\top}\varphi$ with $MM^{\!\top}\varphi=M\bar x$, and
show $MM^{\!\top}=A_{\rm FP}A_{\rm FP}^{\!\top}+\delta^{-2}I$.
(b) Hence the correction is: $(\rho,m)$ get $A_{\rm FP}^{\!\top}\varphi$ exactly as in the
balanced code, and $s$ gets $s=\bar s+\varphi/\delta$. Check the sign.
(c) The balanced code solves $A_{\rm FP}A_{\rm FP}^{\!\top}\varphi=f$ by a DCT in space, which
decouples the spatial modes $k$, followed by one tridiagonal (CN) or ETD system in time per
mode. Explain why adding $\delta^{-2}I$ adds $\delta^{-2}$ to the diagonal of every one of
those time systems and changes nothing else.
(d) The DC mode. For $k=0$ (spatially constant mode) the balanced operator is the Neumann
Laplacian in time, which is singular: its null vector is the constant in $t$, and solvability
requires the residual's DC mode to sum to zero over $t$, which is mass conservation. Show
that $+\delta^{-2}$ makes it regular, and interpret: the total mass at each time slice is now
a free variable, paid for through $s$.
\end{ex}

\begin{codebox}{projection\_unbalanced.py}
\code{inv\_d2 = 1/delta**2} added to \code{D} (interior rows) and to \code{D[0]},
\code{D[-1]}; \code{dc\_den = problem.lambda\_t[:,0] + inv\_d2}; the residual \code{f} has
\code{- s/delta}; the DC mode is solved by a DCT in $t$ divided by \code{dc\_den};
\code{s\_out = s + phi/delta}. Match (a)--(d). Confirm the per-mode off-diagonal \code{e}
is untouched.
\end{codebox}

\begin{ex}{Why the LADMM loop itself does not change.}
The coupling operator $A$ is the interpolation on $(\rho,m)$ and the \emph{identity} on $s$
(\code{A\_fn}, \code{At\_fn} in \code{pipeline\_unbalanced.py}). Why does this keep
$\|A\|\le1$, and why does that matter for the step sizes $\gamma,\tau$ and the stopping rule?
Where does $s$ live on the grid, and why there? (Hint: where does the FP residual live?)
\end{ex}

\begin{ex}{Cost per iteration.}
Count arrays and kernels. In $d$ dimensions the balanced state has $1+d$ fields; the
unbalanced one has $2+d$. The projection does the same DCT and the same batched time solve,
plus one cheap DC-mode solve. Estimate the relative increase in memory and in elementwise
work in 3D. Why can the iteration count go \emph{down} (T1: 53 iterations), and in which
regime does it go up (T3)?
\end{ex}

\section*{Part C. Code audit checklist (1D)}

\begin{enumerate}
  \item \code{StateU}: third field \code{s} of shape \code{(nt, nx)} in both $x$ and $y$;
        \code{norm} includes it. \code{pipeline\_unbalanced.norm\_fn} weights it by
        $dt\,dx$ like $\rho,m$: is $s$ a density-scale quantity? ($s=\delta\zeta$ has units of
        $\rho$ times length/time; on the unit box with $\delta\sim1$ it is comparable to $m$.)
  \item Initial guess: $\rho$ linear in $t$ between $\rho_0,\rho_1$, $m=s=0$. With unequal
        masses the linear interpolant violates the balanced constraint by a constant-in-$x$
        amount; which variable absorbs it on the first projection?
  \item Dual \code{solve\_y}: \code{prox\_ke\_reaction\_cc(z\_hat - delta\_mult/gamma, 1/gamma)}.
        The name \code{delta\_mult} is the LADMM multiplier, not the WFR $\delta$. Rename
        mentally.
  \item Diagnostics returned: \code{ke}, \code{reaction}, \code{mass\_t}. Test T6: the
        discrete mass balance $\tfrac{d}{dt}\int\rho=\int s/\delta$ holds to $10^{-13}$.
        Reproduce it from \code{mass\_t} and \code{s\_stag}.
  \item Run \code{scripts/validate\_unbalanced.py} once and read every printed line against
        the table in \code{unbalanced\_1d.md}.
\end{enumerate}

\section*{Part D. What the symbols mean in a rat-brain tracer scan}

\begin{physbox}{{The experiment (Chen et al.\ 2024, Methods; details not stated there are flagged)}}
A 3-month-old healthy rat, anaesthetized. The MRI contrast agent gadoteric acid (Gd-DOTA) is
infused into the cerebrospinal fluid, in this lab's protocol via a cannula in the cisterna
magna at the base of the skull. T1-weighted 3D images are acquired every 5 minutes for 140
minutes, 29 frames. Gadolinium shortens $T_1$, so the signal rises where tracer is; the
authors convert each frame to \emph{percent signal change from the pre-injection baseline}
and treat that as proportional to tracer concentration. The analysis uses every other frame,
15 inputs 10 minutes apart (repository file names \code{images\_7\_tind\_1} to
\code{images\_35\_tind\_15}), restricted to a brain mask, on a $56\times106\times51$ grid,
smoothed (\code{smooth=1} in the run folder). Field strength, voxel size and the mask
construction are not stated in the paper; read them off the data files. Successive pairs are
solved in a chain, the interpolated end state of one pair seeding the next.
\end{physbox}

\begin{ex}{Dictionary.}
Fill in what each of our symbols is in the scan, then compare with the answers.
\begin{center}\small
\begin{tabular}{@{}lll@{}}
\toprule
ours & in the scan & caution \\
\midrule
$\rho(t,x)$ & \rule{3.2cm}{0.3pt} & linear in Gd only at low concentration \\
$\int\rho\,dx$ & \rule{3.2cm}{0.3pt} & arbitrary units; only the \emph{ratio} $M_1/M_0$ is physical \\
$v=m/\rho$ & \rule{3.2cm}{0.3pt} & \\
$\eps$ & \rule{3.2cm}{0.3pt} & see Exercise~\ref{ex:eps} \\
$\zeta=\rho g$, $g>0$ & \rule{3.2cm}{0.3pt} & \\
$g<0$ & \rule{3.2cm}{0.3pt} & \\
$\delta$ & \rule{3.2cm}{0.3pt} & \\
$t\in[0,1]$ & \rule{3.2cm}{0.3pt} & \\
\bottomrule
\end{tabular}
\end{center}
\end{ex}

\begin{ex}{urOMT to ours: the source.}
Their model: $\partial_t\rho+\dv(\rho v)=\sigma\Delta\rho+\chi\rho r$ with cost
$\int\!\!\int\rho|v|^2+\alpha\chi\rho r^2$, no $\tfrac12$, $\chi\in\{0,1\}$ an indicator
of where sources are allowed, $\chi\equiv1$ for the rat. Match $\zeta=\chi\rho r$, $g=r$ where
$\chi=1$, and show that dividing their integrand by $2$ gives ours with
\[
  \delta^2=\alpha\qquad\text{(in their units, voxels}^2\text{)}.
\]
Now nondimensionalize. Their grid has $\Delta x=1$ voxel and a pair of frames spans
$T=m\,\Delta t=10\times0.4=4$ time units. Rescaling $x\to x/L$, $t\to t/T$ with $L$ the box
side in voxels, show $v\to vT/L$, $g\to gT$, and hence
\[
  \delta_{\rm ours}=\sqrt{\alpha}/L .
\]
Evaluate for the rat ($\alpha=10^4$, $L=106$, the longest axis) and for the spheres
($\alpha=9000$, $L=50$). Which regime of Exercise~3 is this? What does that say about how
the optimizer will use the source: to teleport mass across the brain, or to adjust mass
locally where $\rho$ is already large?
\end{ex}

\begin{ex}{urOMT to ours: the diffusion, and is it physical?}
\label{ex:eps}
Their $\sigma=0.002$ voxel$^2$ per time unit. Show $\eps_{\rm ours}=\sigma T/L^2$ and
evaluate for $L=106$ and $L=50$. Then put units back. Assume a voxel of $0.3$~mm (typical for
this lab's $9.4$~T rat protocol; confirm from the data) and one time unit $=10\,\mathrm{min}/4=150$~s.
Convert $\sigma$ to mm$^2$/s. Compare with the free diffusivity of Gd-DOTA in water at body
temperature, about $4\cdot10^{-4}$~mm$^2$/s, reduced by tissue tortuosity to roughly
$1.5\cdot10^{-4}$~mm$^2$/s. By what factor is their $\sigma$ below the physical value? Finally
compute the \emph{physical} $\eps$ for one 10-minute pair on the $L=106$ crop. Is it inside the
range the 3D suite has already run ($10^{-3}$) and the 2D suite ($10^{-3}$ to $10^{-1}$)?
Which of the two values would you run, and why both?
\end{ex}

\begin{ex}{Where mass comes and goes.}
Over a 10-minute pair the total tracer in the field of view changes because (a) tracer is
still arriving from the infusion site, (b) tracer leaves along perineural and meningeal
lymphatic routes out of the field of view and is cleared to blood, (c) the signal-to-
concentration map saturates at high concentration near the injection site, so early frames
\emph{under}count mass there, (d) noise and partial volume. Which of these is a genuine
source $\zeta$ in the sense of the model, and which is a measurement artefact the model will
nevertheless absorb into $\zeta$? Given the Fisher--Rao cost $\zeta^2/\rho$, where will the
optimizer prefer to place mass gain and loss, and does that match the paper's observation of
early $r>0$ in CSF spaces and later $r<0$ along the skull base?
\end{ex}

\begin{ex}{Hard endpoint versus their free endpoint.}
They solve with $\rho(0)=\rho_0^{\rm img}$ fixed and a penalty $\beta\|\rho(T)-\rho_1^{\rm img}\|^2$,
with $\beta=50$ for the rat (versus $5000$ for the clean spheres). We impose
$\rho(1)=\rho_1$ exactly. (a) With a small $\beta$ their end state may differ visibly from
the image; with ours the image's noise is imposed exactly and must be produced by $v$ and
$\zeta$. Which quantity absorbs pixel noise in our formulation, and at what cost scale (use
$\zeta^2/\rho$ with $\zeta\sim$ noise, $\rho\sim$ local level)? (b) Argue that smoothing
$\rho_0,\rho_1$ with the same kernel is the cheap substitute for a soft endpoint. (c) What
would a soft endpoint change in our LADMM: the prox, the projection, or the coupling? (Only
the terminal row of $A_{\rm FP}$: $\rho_1$ becomes an unknown with a quadratic penalty, so
the time systems get a modified last row. It is not part of today.)
\end{ex}

\section*{Part E. Setting up $\rho_0$ and $\rho_1$}

\begin{enumerate}
  \item \textbf{Pick the pair.} Two consecutive inputs, 10 minutes apart. Start where the
        mass change is clearly nonzero and the geometry is simple: an early pair (tracer
        entering, $M_1>M_0$) and a late pair (clearance, $M_1<M_0$). Record $M_0,M_1$ in the
        data's units before doing anything else.
  \item \textbf{Values.} Percent signal change, masked to the brain. Clip negatives to zero
        (noise below baseline), then add a floor of $10^{-3}$ of the peak inside the mask and a
        smaller floor outside it, so the density is strictly positive everywhere on the box.
        The prox tolerates $\rho\to0$, but the hard endpoint with $\rho_1=0$ where $\rho_0>0$
        forces that mass to be destroyed, which is a real statement; make it only inside the
        mask.
  \item \textbf{Smoothing.} One Gaussian pass, same kernel on both frames, width about one
        voxel; this is what their \code{smooth=1} does. It is also the substitute for their
        soft endpoint (Exercise~15).
  \item \textbf{Geometry.} The crop is $56\times106\times51$ voxels, isotropic. Our 3D setup
        takes separate $n_x,n_y,n_z$ but one side length $L$, so unequal counts would stretch
        the voxels. Pad the volume with the outside floor to a cube of $106^3$ voxels, or
        downsample by two to $28\times53\times26$ and pad to $56^3$ (or $64^3$) for the first
        run. Padding with near-zero density costs nothing in transport and keeps the Neumann
        walls away from the tracer. Their boundary is also a Neumann box, not the brain surface.
  \item \textbf{Scale.} Divide both frames by $M_0$ so that $\sum\rho_0\,dV=1$ and
        $\sum\rho_1\,dV=M_1/M_0$. The mass ratio is the physical content; the overall scale
        sets the density peak, which sets $\gamma$ (the balanced 2D/3D runs want
        $\gamma\approx0.1$ at peaks of order $10$--$60$). Do \emph{not} renormalize $\rho_1$ to
        unit mass: that is exactly what the balanced code does and what the unbalanced model
        exists to avoid.
  \item \textbf{Parameters.} $\delta=\sqrt\alpha/L$ and $\eps=\sigma T/L^2$ from Exercises~12
        and~13 with $L$ the padded cube side in voxels, so the first run reproduces their
        setting; then $\eps_{\rm phys}$ as a second run.
  \item \textbf{Conversion back.} Velocity in voxels/min $=v\,L/T_{\rm min}$ with
        $T_{\rm min}=10$; multiply by the voxel size for mm/min. Growth rate in 1/min $=g/T_{\rm min}$;
        their $r$-map is our $s/(\delta\rho)$ averaged over $t$.
\end{enumerate}

\begin{ex}{Sanity numbers before the run.}
From $\rho_0,\rho_1$ alone, compute: $M_1/M_0$; the Hellinger cost $2\delta^2\!\int(\sqrt{\rho_1}-\sqrt{\rho_0})^2$
(the cost of pure reaction, an upper bound on the optimum); the balanced $W_2$ cost of the
renormalized pair from a cheap Sinkhorn or from the solver's balanced mode (roughly the
transport part). Their ratio tells you in advance whether this pair is reaction-dominated or
transport-dominated at the chosen $\delta$. Also compute the mass fraction within two voxels
of the box walls after padding; it should be below $10^{-4}$.
\end{ex}

\section*{Part F. Porting to 3D and running today}

The same three edits as in 1D, in \code{3d/setup/}; nothing in the DCT structure, the jitted
loop or the mixed-precision path changes.
\begin{enumerate}
  \item \code{state\_3d.py}: add field \code{s} of shape \code{(nt, nx, ny, nz)} to the state
        (both copies), include it in the norm and the arithmetic.
  \item \code{prox\_3d.py}: \code{q2 = mx**2 + my**2 + mz**2 + s**2} (on the cell-centre copy),
        shrink \code{s} by the same factor as the momenta.
  \item \code{projection\_expsemi\_3d.py} (and \code{projection\_3d.py} if the CN path is
        kept): add \code{1/delta**2} to the diagonal of the per-mode time system in
        \code{precomp}, replace the DC-mode special case by the regular solve with
        \code{lambda\_t + 1/delta**2}, subtract \code{s/delta} in the residual, add
        \code{phi/delta} to \code{s} in the output. The ETD coefficients are unchanged.
  \item \code{pipeline\_3d.py}: \code{A\_fn}/\code{At\_fn} identity on \code{s}; initial
        \code{s = 0}; return \code{ke}, \code{reaction}, \code{mass\_t}.
  \item \textbf{Tests before data}, in this order, at $32^3$ then $48^3$: T1 uniform
        $\to2\times$uniform with $\eps=0.01$, $\delta=0.3$ (cost $2\delta^2(\sqrt2-1)^2$, exact for
        any $\eps$, converges in tens of iterations); T4 equal-mass Gaussians with
        $\delta=0.5,1,2,4$ against the balanced solver (difference $\propto1/\delta^2$); T6
        mass balance to machine precision. These three take minutes and catch every sign.
  \item \textbf{Spheres.} Their synthetic case is public: five frames of a translating,
        diffusing Gaussian on $50^3$ with mass factor $(1+a_i\chi_i)$, $a=(0,0.1,0.2,0.1,0)$,
        $\chi_i$ a radius-1.5 ball around the centre. Run pair $0\to2$ with $\delta=\sqrt{9000}/50$,
        $\eps=3.2\cdot10^{-6}$ (their setting; the solver has run $10^{-3}$ at $96^3$ and the
        1D code handles $\eps=0$, so expect no difficulty) and compare the mass curve
        \code{mass\_t} with $(1+a_i)$ and the location of $\zeta>0$ with the ball. Note their
        blur per frame, std $(i+1)\sqrt{0.2}$ voxels, is far more than $2\sigma T$ accounts for;
        the unbalanced solver will spend some of its source on that mismatch.
  \item \textbf{Rat pair.} One early and one late pair as in Part~E, at the padded $56^3$
        first. Report $M_1/M_0$, kinetic and reaction energies, the mass curve, a
        time-averaged $g$ map alongside their Fig.\ (early $r>0$ in CSF, later $r<0$ at the
        skull base), and a speed map. Then $106^3$ in \code{mixedc} precision if time allows.
  \item \textbf{Write-up.} State the two parameter translations with their numbers, the hard
        versus soft endpoint difference, the padding, and that $\eps$ at their value is a
        regularizer while $\eps_{\rm phys}$ is $10^2$ larger; give the mass curves and energy
        split for both $\eps$ values.
\end{enumerate}

\appendix
\section{Answers and hints}
\label{sec:answers}

\paragraph{Ex.\ 1.} $v$ has units length/time, $g$ has units 1/time, so $\delta$ is a length.
The metric $ds^2=|v|^2+\delta^2g^2$ needs a length to compare moving mass with changing it.

\paragraph{Ex.\ 2.} With $\rho=\sigma^2$, $(\partial_t\rho)^2/\rho=4(\partial_t\sigma)^2$, and
$\int_0^1(\partial_t\sigma)^2dt$ with fixed endpoints is minimized by the straight line. Uniform
$\rho_0$ stays uniform, so $\Delta\rho=0$ and the diffusion term never acts.

\paragraph{Ex.\ 3.} Naive crossover $\tfrac12d^2=4\delta^2$, $d=2\sqrt2\,\delta$. The exact
geodesic shrinks mass while moving it and recreates it on arrival, with cost
$\propto\delta^2(1-\cos(d/2\delta))$, finite only for $d<\pi\delta$; beyond that the geodesic
is pure teleport. In T3$'$ the optimizer does exactly that: destroys part of the mass at the
start, moves the rest, and regrows it at the target, so the total mass dips in the middle.

\paragraph{Ex.\ 4.} Lagrangian $\iint\tfrac12\rho(|v|^2+\delta^2g^2)-\phi(\partial_t\rho+\dv(\rho v)-\eps\Delta\rho-\rho g)$;
vary $v$: $\rho v=\rho\nabla\phi$; vary $g$: $\delta^2\rho g=\rho\phi$. So
$\phi=\delta^2\zeta/\rho$ and $v=\nabla\phi$. As $\delta\to\infty$ with $\phi$ fixed,
$\zeta=\rho\phi/\delta^2\to0$: the source is switched off and the constraint becomes the
balanced one, with error $O(1/\delta^2)$.

\paragraph{Ex.\ 5.} (a) $m/\rho+(m-\bar m)/\sigma=0\Rightarrow m=\bar m\rho/(\rho+\sigma)$, same
for $s$. (b) Substituting, the $m,s$ part of the objective becomes $q^2/(2(\rho+\sigma))$;
differentiating $q^2/2(\rho+\sigma)+(\rho-\bar\rho)^2/2\sigma$ in $\rho$ and clearing
denominators gives the cubic. (c) It is the perspective of $\tfrac12|(m,s)|^2$, jointly
convex in $(\rho,m,s)$ because the perspective of a convex function is convex. (d)
$\rho\,h(\zeta/\rho)$ with $h(g)=g\log g-g+1$ or $|g|$: the $s$-stationarity is no longer
linear, so a scalar Newton iteration per cell replaces Cardano.

\paragraph{Ex.\ 6.} (a) $MM^{\!\top}=A_{\rm FP}A_{\rm FP}^{\!\top}+(1/\delta)^2I$ because the
$s$ block of $M$ is $-I/\delta$. (b) $x=\bar x-M^{\!\top}\varphi$: the $s$ component of
$M^{\!\top}\varphi$ is $-\varphi/\delta$, so $s=\bar s+\varphi/\delta$. (c) The DCT
diagonalizes the spatial part; $\delta^{-2}I$ is already diagonal in every basis and
commutes with everything, so it lands on the diagonal of each mode's time system. (d) The
Neumann time Laplacian has eigenvalues $\lambda_t\ge0$ with one zero; $\lambda_t+\delta^{-2}>0$.
The DC component of the residual is the mass defect per time slice, and it is now absorbed
by $s$ instead of being required to vanish.

\paragraph{Ex.\ 7.} $A$ is block diagonal with an interpolation block of norm $\le1$ and an
identity block of norm $1$. LADMM convergence needs $\tau\ge\gamma\|A\|^2$; with $\|A\|\le1$
the balanced choice $\tau=1.01\gamma$ still works. $s$ sits on cell centres in $t$ and $x$
because that is where the FP residual is evaluated, so the constraint couples them without
interpolation.

\paragraph{Ex.\ 8.} 3D: $5$ fields instead of $4$, so about $+25\%$ memory for the state and
the same fraction in the prox and residual; the projection's DCT and time solves are
unchanged in count. Iterations fall when the DC mode was the slow mode (it is now regular);
they rise when $\delta$ is small enough that teleporting is optimal, since the problem then
has flat directions in where exactly the mass dies and is reborn.

\paragraph{Ex.\ 11.} $\rho$: tracer concentration proxy (percent signal change), masked.
$\int\rho$: total tracer in the field of view. $v$: bulk cerebrospinal-fluid flow along
perivascular spaces, the quantity their pathlines draw. $\eps$: tracer diffusivity in
tissue (Exercise~13 says theirs is far below physical). $g>0$: tracer arriving, from the
infusion or from outside the field of view, or apparent gain from signal saturation
relaxing. $g<0$: clearance to lymphatics and blood, tracer leaving the field of view.
$\delta$: the distance beyond which it is cheaper to clear and re-inject than to move;
$\sqrt\alpha=100$ voxels for the rat. $t\in[0,1]$: the 10 minutes between the two frames.

\paragraph{Ex.\ 12.} $\tfrac12(\rho|v|^2+\alpha\rho r^2)=\tfrac12\rho(|v|^2+\delta^2g^2)$ with
$\delta^2=\alpha$. Under $x'=x/L$, $t'=t/T$: $v'=vT/L$, $g'=gT$, so
$|v|^2+\alpha g^2=(L/T)^2(|v'|^2+(\alpha/L^2)g'^2)$ and $\delta'=\sqrt\alpha/L$. Rat:
$100/106=0.94$; spheres: $94.9/50=1.90$. Both are about one box length, so $\pi\delta$
exceeds the brain: the source will never teleport mass across the domain; it adjusts the
total locally, preferring places where $\rho$ is large because the cost is $\zeta^2/\rho$.

\paragraph{Ex.\ 13.} $\eps=\sigma T/L^2$: $0.002\cdot4/106^2=7.1\cdot10^{-7}$ and
$0.002\cdot4/50^2=3.2\cdot10^{-6}$. Units: $0.002\,\text{voxel}^2/150\,\text{s}\times0.09\,\text{mm}^2/\text{voxel}^2=1.2\cdot10^{-6}$~mm$^2$/s,
about $100\times$ below the tissue diffusivity of Gd-DOTA. Physical:
$\eps=D\,T/L^2=1.5\cdot10^{-4}\cdot600/(106\cdot0.3)^2\approx9\cdot10^{-5}$, so $\sim10^{-4}$,
inside the ranges already run. Run both: theirs to compare with their maps, the physical
one because it is the value the paper's own ``$\eps$ is tracer diffusivity'' claim implies.

\paragraph{Ex.\ 14.} (a) and (b) are genuine sources at the scale of the model (mass crossing
the field-of-view boundary or the blood--CSF barrier); (c) and (d) are artefacts that the
model will still book as $\zeta$. The Fisher--Rao cost favours gain and loss where $\rho$ is
large, i.e.\ in the bright CSF spaces, consistent with their maps.

\paragraph{Ex.\ 15.} (a) $\zeta$ absorbs it: pixel noise of amplitude $\eta$ on a level $\rho$
costs $\sim\delta^2\eta^2/\rho$ per cell per unit time, cheap where $\rho$ is large and
expensive in dim tissue, so noise in dim regions is instead transported in from neighbours.
(b) Smoothing both frames with one kernel removes the component that a soft endpoint would
have ignored, without changing the solver. (c) Only the terminal row of the constraint.

\end{document}
